\section{置換で形を変えてから部分積分する}\label{節:置換と部分積分}

積分法は，問題ごとに一つだけ選ぶものではない。置換によって合成関数をほどくと，今度は二つの関数の積が現れ，部分積分が使えることがある。前章でも，根号を変数に換えた後で割り算を使った。同じように，途中の式に合わせて次の道具を選べばよい。

\begin{example}{対数を角に持つ三角関数}
$x>0$ において，$\int\sin(\log x)\dd x$ を求めよ。
\end{example}

角の $\log x$ を一つの変数 $t$ にしたい。しかし，もとの積分には $\frac1x$ がないので，$\sin t\dd t$ と置き換えることはできない。$t=\log x$ を $x=\e^t$ と解き直すと，
\[
 \dd x=\e^t\dd t
\]
である。この微分の係数も含めて換えると，
\[
 \int\sin(\log x)\dd x
 =\int\e^t\sin t\dd t
\]
となる。これは直前の例題で求めた形だから，
\[
\begin{aligned}
 \int\sin(\log x)\dd x
 &=\frac12\e^t(\sin t-\cos t)+C\\
 &=\frac x2\{\sin(\log x)-\cos(\log x)\}+C.
\end{aligned}
\]
$\log x$ は $x>0$ で狭義単調増加し，逆向きの対応 $x=\e^t$ も一つに決まっている。置換後に $\e^t$ が現れたのは，新しい問題を余計に作ったためではなく，$\dd x$ を正しく書き換えたためである。

検算では，外にある $x$ と，内側の $\log x$ の両方を微分する。
\[
\begin{aligned}
 &\frac{\dd}{\dd x}
 \left[\frac x2\{\sin(\log x)-\cos(\log x)\}\right]\\
 &\qquad=\frac12\{\sin(\log x)-\cos(\log x)\}
   +\frac x2\cdot\frac1x
     \{\cos(\log x)+\sin(\log x)\}\\
 &\qquad=\sin(\log x).
\end{aligned}
\]
$t=\log x$ と置くだけでは終わらず，置換後の積を部分積分で求め，最後に元の変数へ戻した。この順序が，それぞれの手法の役割を分けている。

